% Separate post-CP29-build candidate; no change to the frozen current gate.
\subsection{Actual terminal matrices for a reducible first-wild comparison}

Let \(p\ge5\), \(q=p^2\), \(N=q-1\), and let \(K\) be complete
discretely valued of mixed characteristic \((0,p)\), with
algebraically closed residue field, \(\mu_p\subset K\), and
\(e=e_K=p(p-1)\). Put \(s=(p+1)e\) and take the common
tame good field \(L/K\) of degree \(N\). All groups below
are whole potentially supersingular height-two groups; the
comparisons are of their actual generic \(p^3\)-torsion.

\begin{lemma}[Unequal slopes at the first wild index]
\label{lem:B-RD82-Newton}
\status{Candidate; full proof supplied, independent audits pending.}
If the first-level Newton polygon has two slopes, there is a
canonical high line \(H\) and an integer \(B\) such that
\[
 \begin{gathered}
 0<B<pe,\qquad \alpha=s-B,\qquad \beta=B/p,\\
 v_L(T(Q))=\alpha\ (0\ne Q\in H),\qquad
 v_L(T(Q))=\beta\ (Q\notin H).
 \end{gathered}                                      \tag{B-RD82.1}
\]
The high line is \(L\)-rational. If \(c\) is the finite
special tangent exponent, its actual tame character is
\(\zeta^{c-B}\), where \(\zeta\) is the generator chosen
on a tame uniformizer of \(L\). A primitive third-level
point with \(p^2P\notin H\) has
\[
                    v_L(T(P))=B/p^5.                 \tag{B-RD82.2}
\]
The residual division field over \(L\) has degree one or
\(p\). If it has degree \(p\), its unique wild fixed line
is \(H\).
\end{lemma}
\begin{proof}
Here the coefficient proof of B60 and B69 is reapplied at the
present index, rather than importing their prime-to-\(p\)
field scope. In the actual universal height-two coordinate,
the only possible negative Newton vertices of \([p]\) are
\((1,eN),(p,u_0),(q,0)\). The coefficient estimates are
those proved in Lemma~\ref{lem:B-RD49-point}. The unequal
polygon has \(p-1\) points at valuation
\((eN-u_0)/(p-1)\) and \(q-p\) at valuation
\(u_0/(p(p-1))\). The first points and zero form a
unique inertia-stable line. Its character is
\(\mathbf F_p^\times\)-valued and hence is killed by \(L\).
Its coordinate valuation is therefore integral. As
\(eN=(p-1)s\), this gives \(u_0=(p-1)B\), with the
range and values (B-RD82.1).

Use an integral Honda framing and tame averaging as in
Lemma~\ref{lem:B-RD67-normalization}. The leading action
on a high point is \(\zeta^{c+\alpha}=\zeta^{c-B}\),
because \(s\) is a multiple of \(N\). This is an actual
line character, and therefore remains identified by any
torsion isomorphism identifying the two high lines.

At a low target of valuation \(B/p\), the inverse polygon
has one segment from \((0,B/p)\) to \((q,0)\).
Indeed \(u_0=(p-1)B>B/p\), the linear coefficient has
valuation \(eN>B/p\), and the intervening coefficients
have the same lower bounds as in B60. Every inverse has
valuation \(B/p^3\). The still smaller target of that
valuation gives the same one-segment argument and
valuation \(B/p^5\), proving (B-RD82.2).

The residual high line and residual determinant are trivial
over \(L\): the latter is cyclotomic, with \(\mu_p\subset L\).
Thus the residual image over \(L\) lies in the upper
unipotent group of order \(p\). If nontrivial, its unique
fixed line is the already invariant high line.
\end{proof}

\begin{lemma}[An actual terminal plane forced by a prime point valuation]
\label{lem:B-RD82-terminal}
\status{Candidate; full proof supplied, independent audits pending.}
Suppose one such group has two slopes with \(B=pb\),
\(p\nmid b\), and its full third-level division field
\(F=L(M[p^3])\) has \([F:L]=p^4\). In matrix coordinates
with \(H=\mathbf F_p e_1\), its terminal matrix space
\[
 V_2=\{C\in M_2(\mathbf F_p):1+p^2C\in
                \operatorname{Gal}(F/L)\}
\]
has dimension two. It contains \(E_{12}\) and a matrix
whose induced coefficient on the quotient by \(H\) is
nonzero. Every element of \(V_2\) is upper triangular,
and every element with zero lower diagonal coefficient
is a multiple of \(E_{12}\).
\end{lemma}
\begin{proof}
Write \(W=\operatorname{Gal}(F/L)\). Its residual image
has order \(p^\epsilon\), with \(\epsilon\in\{0,1\}\).
By (B-RD82.2), every primitive point whose residual class
is outside \(H\) has integral \(F\)-valuation \(b\),
which is prime to \(p\).

For a nonidentity terminal element \(g=1+p^2C\), the
fixed-coefficient uniformizer argument of
Lemma~\ref{lem:B-RD60-break} gives
\[
             v_F(gT(P)-T(P))=b+b(g)                 \tag{B-RD82.3}
\]
for every such primitive point. In particular \(Cv\ne0\)
for every residual vector \(v\notin H\). Moreover all
these images \(Cv\) have the same high or low status:
their displacement valuations are either \(p^4\alpha\)
or \(p^4b\), and (B-RD82.3) is independent of \(v\).
Formal subtraction and coordinate subtraction have the
same valuation, since their ratio is an integral unit.

Write the lower row of \(C\) as \((c,d)\). If \(c\ne0\),
the vectors \((x,1)\) include one with zero lower image
and one with nonzero lower image. The first image cannot
be zero, and hence would be high, while the other is low.
This contradicts the common status just proved. Thus
\(c=0\). If \(d=0\), the upper row \((a,b_0)\) cannot
have \(a\ne0\), since it would then kill one vector
\((x,1)\). Consequently \(C=b_0E_{12}\).
The kernel of the linear map \(V_2\to\mathbf F_p\)
taking the lower diagonal coefficient is therefore
contained in \(\mathbf F_pE_{12}\), so \(t=\dim V_2\le2\).

Let \(V_1\) be the image of
\(\ker(W\to\GL_2(\mathbf F_p))\) under the first
congruence coefficient, and put \(r_1=\dim V_1\).
The actual matrix identity
\((1+pA+p^2B)^p=1+p^2A\pmod{p^3}\) gives
\(V_1\subset V_2\). The exact group-order filtration
therefore reads
\[
             4=\epsilon+r_1+t,
                    \qquad r_1\le t\le2.             \tag{B-RD82.4}
\]
It forces \(t=2\), with \(r_1=2\) for \(\epsilon=0\)
and \(r_1=1\) for \(\epsilon=1\). The lower-diagonal
map is now surjective and its one-dimensional kernel
is exactly \(\mathbf F_pE_{12}\). This proves the
assertion using actual terminal elements; no bound on
integral orbit sizes is inferred from residual evaluation ranks.
\end{proof}

\begin{proposition}[Recovering the high line from two terminal breaks]
\label{prop:B-RD82-recovery}
\status{Candidate; full proof supplied, independent audits pending.}
Let two groups with two-slope first-level polygons have
isomorphic actual \(p^3\)-torsion and common division
field \(F/L\) of degree \(p^4\). If one parameter is
\(B_1=pb_1\) with \(p\nmid b_1\), then their high lines
agree under the actual isomorphism, \(B_1=B_2\), and
their finite special characters agree.
\end{proposition}
\begin{proof}
Choose a residual vector outside both high lines. Its
primitive lifts in model 2 have valuation \(B_2/p^5\)
over \(L\). Integrality in the actual common field of
degree \(p^4\) implies \(p\mid B_2\). Write
\(B_2=pb_2\); both integers satisfy \(0<b_i<e\).
The preceding lemma supplies an actual terminal element
with matrix \(E_{12}\) relative to \(H_1\). For every
vector outside \(H_1\), its displacement in model 1 is
high, so its exact lower break is
\[
                  r_N=p^4(s-pb_1)-b_1.               \tag{B-RD82.5}
\]
If \(H_1\ne H_2\), choose the residual vector outside
both. Its nonzero image under \(E_{12}\) lies in
\(H_1\setminus H_2\), so the corresponding displacement
in model 2 is low. The universal lower bound, which
does not require \(p\nmid b_2\), would give
\[
       p^4b_2\ge b_2+p^4(s-pb_1)-b_1.                \tag{B-RD82.6}
\]
This is impossible: \(s-pb_1>e\), whereas both
\(b_1,b_2<e\). More explicitly the left side minus
the right side is increasing in both integers and,
at \(b_1=b_2=e-1\), is already \(-p^4(p+1)<0\).
Thus \(H_1=H_2\).

Take now an actual terminal matrix having nonzero lower
diagonal coefficient. It sends every residual vector
outside the common line to a low vector in both models.
Its exact break in model 1 is \((p^4-1)b_1\).
The universal inequality in model 2 implies
\[
        (p^4-1)(b_2-b_1)\ge0.                        \tag{B-RD82.7}
\]
For the nilpotent terminal element (B-RD82.5), both
displacements are high; the same inequality instead gives
\[
                    (p^5+1)(b_1-b_2)\ge0.            \tag{B-RD82.8}
\]
Hence \(b_1=b_2\). The actual high-line characters
\(\zeta^{c_i-B_i}\) agree under the given isomorphism.
Equality of \(B_i\) therefore identifies \(c_1=c_2\)
modulo \(N\), proving the special-character assertion.
\end{proof}

\begin{theorem}[A lower image bound also for nontrivial residual wild inertia]
\label{thm:B-RD82}
\status{Candidate; full proof supplied, independent audits pending.}
At the stated first wild index, an actual \(p^3\)-congruent
pair with nontrivial residual wild inertia and different
finite special characters has full wild image of order
\(p^r\), with \(5\le r\le9\). In particular its wild
image has order at least \(3125\) when \(p=5,e=20\).
\end{theorem}
\begin{proof}
Their residual wild image over \(L\) has order \(p\).
The remaining two congruence kernels have dimensions at
most four each, giving \(r\le9\). A one-slope first-level
polygon would have all nonzero coordinates at valuation
\(e\), and its injective leading-coordinate map over \(L\)
would make the residual wild action trivial. Thus both
polygons have two slopes and the actual isomorphism
identifies their unique wild fixed high lines.

Suppose \(r\le4\). Equation (B-RD82.2) and integrality
force \(B_i=pb_i\). If one \(b_i\) is prime to \(p\),
the same valuation denominator forces \(r\ge4\), so
\(r=4\); the preceding proposition gives equal characters.

It remains that \(B_i=p^2k_i\). Their range gives
\(1\le k_i\le p-2\). The finite determinant characters
agree by the actual torsion determinant comparison.
If this common character is square,
Proposition~\ref{prop:B-RD77-square} already gives equality.
Otherwise apply a common finite height-one tame twist
to normalize the finite determinant to the primitive
level-one character, exactly as in B50 and B77.
This preserves the actual torsion isomorphism, both
polygons over \(L\), and the wild image. Now
\(c_i\equiv1\pmod{p-1}\). The common high-line character
gives \(c_1-B_1\equiv c_2-B_2\pmod N\), so
\[
                 k_1\equiv k_2\pmod{p-1}.             \tag{B-RD82.9}
\]
The indicated integer interval forces \(k_1=k_2\);
the same line-character equality gives \(c_1=c_2\pmod N\).
This excludes every case with \(r\le4\).
\end{proof}

\begin{remark}[The residual-tame exceptional sector is not removed]
\label{rem:B-RD82-scope}
\status{Candidate scope; no complete wild-index recovery claim.}
The terminal-plane proposition also applies when the residual
representation is tame, provided that its hypotheses on the
two polygons, the common field degree, and one prime-to-\(p\)
integer \(b_i\) hold. The final theorem uses the unique
residual wild fixed line to handle the remaining case with
both \(b_i\) divisible by \(p\). It does not exclude a
residual-tame comparison with two such parameters and
exchanged high lines, nor larger wild images in any branch.
It makes no existence or unrestricted \(p^3\)-recovery
assertion, and the full first-wild index remains in \(\{2,3\}\).
\end{remark}
